% Lösungen Einheit 3 von 5 – Extrem- und Wendepunkte mit Parameter (zusammenbau v0.1)
\einheitenkopf{Einheit 3 von 5 · Extrem- und Wendepunkte mit Parameter}

\erg{14}{a) $f_a'(x) = 3a x^2 - 4x$, $f_a'(1) = 3a - 4$ \quad b) $f_k'(x) = -4k x^2 (x - 6)$: bei $0$ kein Vorzeichenwechsel (Sattelstelle), bei $x = 6$ Wechsel von $+$ nach $-$: Hochpunkt bei $x = 6$. \quad c) $f_a(1) = 2{,}5 + 2a = -0{,}5$, also $a = -1{,}5$; $f''(x) = 3x - 3$, $f''(1) = 0$ und $f'''(1) = 3 \neq 0$: $P$ ist Wendepunkt. \quad d) $f_a'(x) = \frac{2ax}{ax^2 + 4}$; der Nenner ist positiv, also ist $x = 0$ die einzige Nullstelle, und $f_a(0) = \mathrm{ln}(4)$ für jedes $a$; der Zähler $2ax$ wechselt bei $0$ von $-$ nach $+$: Tiefpunkt $(0 \,|\, \mathrm{ln}(4))$. \quad e) $f_a'(x) = \frac{1}{2}(x^2 + 4x + a) \cdot e^{0{,}5x}$; wegen $e^{0{,}5x} > 0$ bleibt $x^2 + 4x + a = 0$ mit $x = -2 \pm \sqrt{4 - a}$; für $a > 4$ ist der Radikand negativ: keine Extrempunkte. \quad f) $K_b'(x) = 3x^2 - 2bx + 12$ hat für $4b^2 - 144 < 0$, also $b < 6$, keine Nullstelle; wegen $K_b'(0) = 12 > 0$ ist $K_b' > 0$ überall. \quad g) $f_k(2k) = \frac{1}{k} \cdot 4k^2 \cdot 4k^2 = 16k^3$, Tiefpunkte $(0 | 0)$ und $(4k | 0)$; Abstand $d = \sqrt{(2k)^2 + (16k^3)^2} = 2k\sqrt{1 + 64k^4}$. \quad h) $f_a'(x) = \frac{3}{a}x^2 + 4x + 3$; genau eine Nullstelle bei Diskriminante $16 - \frac{36}{a} = 0$, also $a = 2{,}25$ (für $a < 0$ ist die Diskriminante stets positiv). Sattelpunkt: $f_a''$ ist an dieser Stelle null und $f_a'''$ ungleich null.}
\erg{15}{a) Die Art wurde ohne hinreichende Bedingung geraten: $f_a''(x) = 6x$, also $f_a''(\sqrt{a}) > 0$ (Tiefpunkt bei $\sqrt{a}$) und $f_a''(-\sqrt{a}) < 0$ (Hochpunkt bei $-\sqrt{a}$). \quad b) $f_a'(x) = 3x^2 + a$ ist für $a > 0$ stets positiv; die Gleichung $f_a'(x) = 0$ hat keine Lösung, also gibt es keine Extremstelle. \quad c) $f_k'(t) = 30(1 - kt)e^{-kt} = 0$ für $t = \frac{1}{k}$ (Maximum); $f_k''(t) = 30k(kt - 2)e^{-kt}$ wechselt bei $t = \frac{2}{k}$ das Vorzeichen: Dort nimmt die Temperatur am stärksten ab.}
